\documentclass[12pt, a4paper]{article}
\usepackage{../../../myconfig} % Gọi file cấu hình từ ngoài gốc vào

% ==========================================
% BẮT ĐẦU NỘI DUNG TÀI LIỆU
% ==========================================
\begin{document}

% Truyền 3 tham số: {Nhãn cam} {Chữ to chính giữa} {Chữ ở góc phải Header}
\titlebox{Chủ đề: Hàm số}{Hàm doanh thu, lợi nhuận, tối ưu chi phí}{Hàm số: Doanh thu lợi nhuận chi phí}

% --- CÂU 5 ---
\questionTrueFalse{0}{1}{Tiệm vải bà Trâu đang tính toán các chi phí để vận hành xưởng và sản xuất \( x \) mét vải trong một tháng. Bà Trâu chia chi phí thành 3 hạng mục chính: chi phí thuê xưởng là 8 triệu đồng, chi phí nguyên liệu để sản xuất một mét vải là 25 nghìn đồng, và chi phí bảo trì máy móc là \( 0,00001x^2 \) triệu đồng. Biết rằng giá vải tiệm yết của tiệm vải bà Trâu là 100 nghìn đồng/mét vải và giả sử toàn bộ số vải sản xuất được thì đều được bán hết.}{
    \textbf{a)} & Tổng chi phí vận hành và sản xuất vải của tiệm vải bà Trâu trong một tháng được tính bằng công thức  
    \( C(x) = 0,00001x^2 + 0,025x + 8 \) (đơn vị: triệu đồng).
    & \textcolor{amberorange}{$\bigcirc$} & \textcolor{amberorange}{$\bigcirc$} \\ \hline
    
    \textbf{b)} & Nếu mỗi tháng tiệm vải bà Trâu sản xuất được 109m vải, khi đó tiệm sẽ bị lỗ.
    & \textcolor{amberorange}{$\bigcirc$} & \textcolor{amberorange}{$\bigcirc$} \\ \hline
    
    \textbf{c)} & Lợi nhuận lớn nhất mà tiệm vải bà Trâu có thể đạt được trong một tháng lớn hơn 132 triệu đồng.
    & \textcolor{amberorange}{$\bigcirc$} & \textcolor{amberorange}{$\bigcirc$} \\ \hline
    
    \textbf{d)} & Giả sử tiệm vải bà Trâu sản xuất số mét vải tiến đến vô cùng, khi đó chi phí trung bình để sản xuất mỗi mét vải sẽ tiến đến 0.
    & \textcolor{amberorange}{$\bigcirc$} & \textcolor{amberorange}{$\bigcirc$} \\
}

% --- CÂU 6 ---
\questionTrueFalse{0}{2}{Anh Huy đẹp trai ngoài việc làm chủ của thương hiệu Huy's Secret, thì anh còn là chủ của thương hiệu ti vi STER. Biết rằng nếu anh Huy định sản xuất \( x \) chiếc ti vi, khi đó tổng chi phí sản xuất là \( C(x) = x^3 - 24x^2 + 350x + 338 \) (triệu đồng). Biết rằng đạo hàm của hàm số \( C(x) \) được gọi là chi phí biên, xấp xỉ với chi phí để sản xuất thêm một đơn vị hàng hóa tiếp theo, tức là đơn vị hàng hóa thứ \( x + 1 \).}{
    \textbf{a)} & Tổng chi phí sản xuất ti vi luôn tăng khi tăng số lượng ti vi cần sản xuất.
    & \textcolor{amberorange}{$\bigcirc$} & \textcolor{amberorange}{$\bigcirc$} \\ \hline
    
    \textbf{b)} & Hàm chi phí biên của sản phẩm ti vi là \( C'(x) = 3x^2 - 48x + 350 \).
    & \textcolor{amberorange}{$\bigcirc$} & \textcolor{amberorange}{$\bigcirc$} \\ \hline
    
    \textbf{c)} & Để chi phí biên nhỏ nhất, anh Huy cần sản xuất đúng \( a \) chiếc ti vi, khi đó \( a \) là số chính phương.
    & \textcolor{amberorange}{$\bigcirc$} & \textcolor{amberorange}{$\bigcirc$} \\ \hline
    
    \textbf{d)} & Để chi phí trung bình nhỏ nhất, anh Huy cần sản xuất đúng \( b \) chiếc ti vi, khi đó \( b \) là số nguyên tố.
    & \textcolor{amberorange}{$\bigcirc$} & \textcolor{amberorange}{$\bigcirc$} \\
}

% --- CÂU 7 ---
\questionTrueFalse{0}{3}{Giá trị của một chiếc máy tính để bàn cao cấp (đơn vị: triệu đồng) được mô hình hóa bởi hàm số  
\(V(N) = \dfrac{3N + 430}{N + 1}\)
trong đó \( N \) là số lần hoạt động của máy tính mỗi ngày. Giả sử rằng anh T là biên tập của team Huy Mathster sẽ sử dụng máy tính  
\(N(t) = \sqrt{t^2 - 10t + 61}\)
lần mỗi ngày, với \( t \) là số tháng kể từ khi mua máy tính.}{
    \textbf{a)} & ANh T sẽ luôn sử dụng máy tính tối thiểu 5 lần mỗi ngày.
    & \textcolor{amberorange}{$\bigcirc$} & \textcolor{amberorange}{$\bigcirc$} \\ \hline
    
    \textbf{b)} & Đạo hàm của hàm số \( V(t) \) là  
    \[V'(t) = \dfrac{-427(t - 5)}{\sqrt{t^2 - 10t + 61} \cdot \left(\sqrt{t^2 - 10t + 61} - 1\right)^2}.\]
    & \textcolor{amberorange}{$\bigcirc$} & \textcolor{amberorange}{$\bigcirc$} \\ \hline
    
    \textbf{c)} & Sau thời gian dài (có thể coi như tiến đến vô cùng), giá trị của máy tính sẽ luôn giảm nhưng không thấp hơn 3 triệu đồng.
    & \textcolor{amberorange}{$\bigcirc$} & \textcolor{amberorange}{$\bigcirc$} \\ \hline
    
    \textbf{d)} & Giá trị của máy tính là cao nhất sau 5 tháng kể từ khi mua máy tính.
    & \textcolor{amberorange}{$\bigcirc$} & \textcolor{amberorange}{$\bigcirc$} \\
}

% --- CÂU 8 ---
\questionTrueFalse{0}{4}{Một công ty kinh doanh trong lĩnh vực vận tải đang vận hành một đội gồm 35 xe chở hàng cùng loại, với lợi nhuận trung bình của mỗi xe là 1 triệu đồng một ngày. Để mở rộng mô hình kinh doanh, công ty dự định bổ sung một số xe chở hàng cùng loại với xe đang vận hành. Công ty đã tiến hành khảo sát và phân tích thị trường, kết quả cho thấy: cứ bổ sung một xe chở hàng cùng loại vào hoạt động thì lợi nhuận trung bình của mỗi xe trong cả đội lại giảm đi 20 nghìn đồng một ngày.}{
    \textbf{a)} & Nếu công ty bổ sung 10 xe chở hàng cùng loại thì lợi nhuận trung bình mỗi ngày của đội xe là 36 triệu đồng.
    & \textcolor{amberorange}{$\bigcirc$} & \textcolor{amberorange}{$\bigcirc$} \\ \hline
    
    \textbf{b)} & Nếu công ty bổ sung \( x \) xe chở hàng cùng loại thì lợi nhuận trung bình mỗi ngày của mỗi xe là \( 1 - 20x \) (triệu đồng).
    & \textcolor{amberorange}{$\bigcirc$} & \textcolor{amberorange}{$\bigcirc$} \\ \hline
    
    \textbf{c)} & Lợi nhuận trung bình mỗi ngày của đội xe được tính bởi công thức \( P(x) = (35 + x)(1 - 2x) \) (triệu đồng).
    & \textcolor{amberorange}{$\bigcirc$} & \textcolor{amberorange}{$\bigcirc$} \\ \hline
    
    \textbf{d)} & Có hai giá trị của \( x \) để lợi nhuận trung bình mỗi ngày của đội xe là lớn nhất.
    & \textcolor{amberorange}{$\bigcirc$} & \textcolor{amberorange}{$\bigcirc$} \\
}


% --- CÂU 19 ---
\questionShortAnswer{0}{5}{Công ty du lịch Trân Đầu muốn mở một tour du lịch bằng du thuyền trên biển, và giám đốc công ty đang tính toán các chi phí cần thiết để vận hành du thuyền. Để tính toán chi phí, giám đốc Trân Đầu đã mô hình hóa lượng dầu diesel tiêu thụ ước tính theo hàm số \( G(x) = \dfrac{250}{x} + \dfrac{x}{10} \) (lít/km), với \( x \) là vận tốc của du thuyền (giả định rằng vận tốc này duy trì không đổi trong suốt quá trình di chuyển).  
Biết rằng giá dầu diesel là 17,5 nghìn đồng/lít. Ngoài ra, công ty cần tuyển thuyền trưởng với mức lương 1,5 triệu đồng mỗi giờ để lái du thuyền này, và quãng đường mà chiếc du thuyền di chuyển trong toàn bộ chuyến du lịch dài 640km. Hỏi thuyền trưởng phải duy trì vận tốc của du thuyền bằng bao nhiêu km/h để tổng chi phí của công ty là thấp nhất (kết quả làm tròn đến hàng phần mười)?}

% --- CÂU 20 ---
\questionShortAnswer{0}{6}{Nhằm quảng bá cho tour du lịch bằng du thuyền trên biển, giám đốc Trân Đầu tìm đến một đơn vị quảng cáo ở nước ngoài. Qua khảo sát, giám đốc thấy rằng nếu công ty chỉ ra \( x \) nghìn đô cho việc quảng cáo hàng năm, thì trung bình mỗi năm công ty sẽ thu hút được khoảng \( M(x) = 2300 + \dfrac{125}{x} - \dfrac{525}{x^2} \) du khách trải nghiệm tour du lịch. Hỏi công ty du lịch Trân Đầu nên đầu tư bao nhiêu đô cho việc quảng cáo hàng năm để thu hút được nhiều khách du lịch nhất?}

% --- CÂU 21 ---
\questionShortAnswer{0}{7}{Chi phí để sản xuất \( x \) chiếc móc khóa in logo Huy Math-ster được mô hình hóa bởi hàm số  
\(C(x) = \dfrac{1}{3}x^2 + 4x + 53\)
(nghìn đồng). Biết rằng quy mô sản xuất móc khóa sẽ thay đổi theo thời gian \( t \), với  
\(x(t) = 0,2t^2 - 0,3t + 30\)
chiếc móc khóa sẽ được sản xuất sau \( t \) giờ kể từ khi bắt đầu sản xuất. Hỏi sau bao nhiêu phút kể từ khi bắt đầu sản xuất thì tổng chi phí sản xuất móc khóa sẽ là nhỏ nhất?}

% --- CÂU 22 ---
\questionShortAnswer{0}{8}{Bạn là chủ của một doanh nghiệp sản xuất ghế công thái học, và anh Huy đang muốn mua một số lượng lớn ghế từ doanh nghiệp của bạn để dành cho đội biên tập, tránh đội biên tập của công ty bị đau lưng. Biết rằng nếu bạn bán mỗi ghế với giá 850 nghìn đồng, anh Huy sẽ mua 15 ghế từ doanh nghiệp của bạn. Nếu bạn giảm giá mỗi ghế là 20 nghìn đồng so với giá gốc, anh Huy sẽ mua thêm 2 ghế từ doanh nghiệp của bạn. Để có được doanh thu lớn nhất, bạn cần bán ghế công thái học cho anh Huy với giá là bao nhiêu nghìn đồng mỗi ghế?}



\end{document}